% Propietario conservado: /Users/ruben/Documents/ChatGPT/jueces y controles/REVISION_ARGUMENTAL_APERTURAS_CIERRES_20260915/delivery/ARTICLE_V_ES_ARGUMENTAL_20260915/manuscrito/sections/30_fock_gibbs.tex
% Recuperación documental para el artículo V; RESULTADO_RECUPERADO.
% Propietario: /Users/ruben/Documents/New project/output/REVISION_EDITORIAL_HMT_MD_20260905/01_FUENTES/INTEGRAL/tree/New project/output/ENSAMBLADO_CAUSAL_RESTAURADO_HMT_MD_20260905/fuente/manuscrito/sucesor_102/deltas_masas/realizacion_equivariante_especies_operadores_masicos_publico.tex
% Líneas de origen: 856--932.
% REV02: transporte functorial; CAR/CCR se construyen en 20.
% Correspondencia editorial: gestion/CONSOLIDACION_PAULI_FOCK_REV02.json.
\section{Transport of Fock completions under refinement}
\label{v:v:rec:fock_gibbs:section:01}

The completions and their operator algebras are constructed in
Section~\ref{v:v:sec:completaciones-algebraicas}. The result of this stage
is their compatibility with refinement morphisms: transport of
the one-particle fiber extends to all occupation sectors.

\subsection{Second quantization of refinement morphisms}
\label{v:v:rec:fock_gibbs:section:02}

Let \(\mathsf{Hol}_{\pm}^{\leq1}\) be the monoidal category of one-particle
fibers, with contractive refinement morphisms preserving
holonomy parity \(\epsilon\in\{+1,-1\}\). It includes, in particular,
unitary intertwiners. The monoidal operation on the fibers is
direct sum; on the completions it induces the tensor product,
graded in the exterior sector. Second quantization uses the
completions of Section~\ref{v:v:rec:completaciones:section:03}:
\[
 \mathfrak F(H,\epsilon)=
 \begin{cases}
 \mathcal F_+(H),
     &\epsilon=+1,\\[2mm]
 \mathcal F_-(H),
     &\epsilon=-1,
 \end{cases}
\]
and, for a contractive intertwiner \(T:H\to H'\),
\[
 \mathfrak F(T)=
 \bigoplus_{n\ge0}\operatorname{Sym}^n(T)
 \quad\text{or}\quad
 \bigoplus_{n\ge0}\bigwedge^n\!T.
\]

\subsection{Functoriality, canonical operators, and transport domains}
\label{v:v:rec:fock_gibbs:section:03}

Let \(\mathcal D_\epsilon(H)\) be the algebraic direct sum of the
finite-particle sectors. In the following identities, \(a_H^\dagger(f)\)
and \(a_H(f)\) denote the creation and annihilation operators of the
chosen completion: the operators \(a\) of the exterior sector or the
operators \(b\) of the symmetric sector in
Section~\ref{v:v:pauli:car-construccion}.

\begin{theorem}[Internal transport of the statistic]
\label{v:v:rec:fock_gibbs:statement:01}
The preceding completion defines a graded monoidal functor on
\(\mathsf{Hol}_{\pm}^{\leq1}\), with bounded operator values and
compatibility with refinement. In the branch \(\epsilon=-1\), creation and
annihilation operators satisfy CAR and occupation operators obey
\(N_s^2=N_s\); in the branch \(\epsilon=+1\), they satisfy CCR. Thus
Pauli exclusion and fermionic and bosonic statistics descend from
holonomy parity once the one-particle representation has been fixed.
For every contractive intertwiner \(T:H\to H'\), \(f\in H\), and
\(g\in H'\), the following hold on \(\mathcal D_\epsilon(H)\):
\begin{align}
 \mathfrak F(T)a_H^\dagger(f)
 &=a_{H'}^\dagger(Tf)\mathfrak F(T),
 \label{v:v:fock:eq:transporte-creacion}\\
 a_{H'}(g)\mathfrak F(T)
 &=\mathfrak F(T)a_H(T^*g).
 \label{v:v:fock:eq:transporte-aniquilacion}
\end{align}
If \(T\) is isometric, the second identity, with \(g=Tf\), also transports
annihilation of \(f\). For \(\|f\|=1\), the occupation operators
then satisfy
\begin{equation}
 N_{H',Tf}\mathfrak F(T)=\mathfrak F(T)N_{H,f},
 \qquad N_{H,f}=a_H^\dagger(f)a_H(f).
 \label{v:v:fock:eq:transporte-ocupacion}
\end{equation}
\end{theorem}

\begin{proof}
Symmetric and exterior powers preserve compositions and identities;
in each degree, their norms are bounded by \(\|T\|^n\). Contractivity
therefore ensures that the direct sum defines a bounded operator. The
decomposition of powers of a direct sum according to their degrees in each
factor gives monoidal compatibility, with the graded exchange sign
in the exterior powers.

The CAR relations and idempotence are
Theorem~\ref{v:v:rec:completaciones:statement:01}; the CCR relations are
\eqref{v:v:pauli:eq:ccr}. To verify their transport, it suffices to act on
symmetric or exterior products of vectors. Applying \(T\) to each factor
commutes with inserting the factor \(f\), which is transformed into \(Tf\);
this proves \eqref{v:v:fock:eq:transporte-creacion}. Annihilation sums the
contractions of each factor and uses
\(\langle g,Tf\rangle=\langle T^*g,f\rangle\); exterior signs and
symmetric normalizations agree on both sides. This yields
\eqref{v:v:fock:eq:transporte-aniquilacion}. If \(T^*T=I\), setting \(g=Tf\)
and composing the two identities proves
\eqref{v:v:fock:eq:transporte-ocupacion}. All these operations preserve
the finite-particle domain.
\end{proof}

For noncontractive morphisms, second quantization is defined first
on \(\mathcal D_\epsilon(H)\); its bounded extension depends on a
uniform bound for the norms over all degrees. Such morphisms lie outside
\(\mathsf{Hol}_{\pm}^{\leq1}\). In the symmetric completion, if
\(\|T\|>1\), the norms \(\|\operatorname{Sym}^nT\|=\|T\|^n\)
grow without bound; in the exterior completion of a finite fiber, only
finitely many degrees are nonzero. The general annihilation
identity retains \(T^*g\); its specialization to the same transported
vector corresponds to the isometry condition stated above.

The distinction is already verified in a one-dimensional fiber. If
\(T=tI\), with \(0<t<1\), and \(e\) is a unit vector, then, on the
one-particle state \(e\), in either completion,
\begin{equation}
 a(te)\mathfrak F(T)e=t^2\Omega,
 \qquad \mathfrak F(T)a(e)e=\Omega,
 \label{v:v:fock:eq:control-no-isometrico}
\end{equation}
where \(\Omega\) is the vacuum. The general identity
\eqref{v:v:fock:eq:transporte-aniquilacion} correctly gives
\(\mathfrak F(T)a(T^*te)e=t^2\Omega\). The factor \(t^2\) identifies
precisely the hypothesis allowing specialization of annihilation
transport and, subsequently, occupation transport.

The result establishes internal algebraic transport. Its identification
with the relativistic spin--statistics theorem retains the additional
conditions of Lorentzian representation and locality specified in
Section~\ref{v:v:rec:completaciones:section:05}.

\clearpage
\section{Thermodynamic realization through Gibbs states and the pressure limit}
\label{v:v:rec:fock_gibbs:section:04}

\subsection{Finite Hamiltonian, partition function, Gibbs state, and cofinal pressure limit}
\label{v:v:rec:fock_gibbs:section:05}

For a compressed Hamiltonian \(H_{X,V}\) over a finite volume or refinement
\(V\), the thermal state is completely defined by
\[
 Z_{X,V}(\beta)=\operatorname{Tr}e^{-\beta H_{X,V}},
 \qquad
 \rho_{X,V,\beta}=Z_{X,V}(\beta)^{-1}e^{-\beta H_{X,V}}.
\]
The thermodynamic limit exists when the pressure
\[
 p_X(\beta)=\lim_{V\nearrow\infty}
 \frac1{\beta |V|}\log Z_{X,V}(\beta)
\]
exists and is independent of the cofinal sequence.

\subsection{Condensation criterion by saturation of the excited-state density}
\label{v:v:rec:fock_gibbs:section:06}

Bosonic condensation is
determined by saturation of the excited-state density; it is not asserted
merely from the existence of a symmetric completion. The problem
is thus expressed by a verifiable mathematical condition on the
Hamiltonian and the limit, rather than by a nominal attribution.
